\chapter{Puntos abiertos y mejoras}\label{ap:abiertos}

Lo que no está cerrado, dicho con claridad. Ninguno impide usar la librería, pero conviene saberlo.

\section{Por comprobar en la maqueta real}
\begin{enumerate}
\item \textbf{Polaridad de \cod{Sensor\_stopper\_abajo}.} La etapa 4 de la cinta avanza cuando este sensor \emph{deja} de estar activo, y el simulador cumple esa condición al instante. El manual lo define como «stopper abajo». Hay que ver con la maqueta, con el stopper subido y bajado, qué valor da el sensor, y si la condición debe ser la contraria.
\item \textbf{Tamaño del bloque, 24 o 28 bytes.} El manual menciona «bloques de 24 bytes» en un apartado y 28 en el resto. La librería usa 28, como el UDT del lector y el Excel.
\item \textbf{Comportamiento del lector real} con tiempos reales (el simulador lo imita con 0,12\,s de proceso).
\item \textbf{Varios palets pegados.} Se comprobó con la estación Out en el simulador antiguo, no en el de la librería.
\end{enumerate}

\section{Por decidir en el grupo}
\begin{enumerate}
\item \textbf{Qué escribe Out exactamente.} El Excel dice que Out escribe el byte 4 (bit 3) y el 18, \emph{y además} pone todo el bloque a cero. Como el cero lo pisa todo, la librería está probada con la escritura completa a cero. Si se prefiere conservar los bytes 4 y 18, la máscara de Out sería \cod{16\#00040010} y no se borraría.
\item \textbf{Calidad «sin evaluar» (byte 5 igual a 0).} Out lo trata como NOK. El mapa no dice nada.
\item \textbf{Qué almacén recibe qué.} En el simulador antiguo la tapa buena iba al almacén derecho; el Excel dice «OK izquierdo, NOK derecho». La aplicación de Out del proyecto actual manda OK a la derecha. Hay que dejarlo escrito en el Excel.
\item \textbf{Job por defecto de la cámara} (en el ejemplo, el 4).
\item \textbf{Bytes 20 a 23}: propuesta de datos de otros grupos, aún sin acordar.
\end{enumerate}

\section{Mejoras para una versión futura}
\begin{enumerate}
\item Leer siempre la unión de la máscara de lectura y la de escritura, para que la regla O1 deje de existir.
\item Un temporizador máximo para la aplicación, con aviso.
\item Quitar la etapa 2 de la cinta, que no se alcanza.
\item Resolver los dos avisos de conversión de enteros de \cod{Est\_FB\_RFID}.
\item Añadir el prefijo \cod{Est\_} a los UDT, para que no choquen con tipos con el mismo nombre en el proyecto de cada uno.
\item Incluir el simulador de pruebas en la distribución de la librería.
\end{enumerate}

\section{Limitaciones de las comprobaciones}
\begin{itemize}
\item La copia de la \emph{tabla de variables} de la plantilla no se ha podido comprobar de manera automática, solo a mano en TIA.
\item Los ejemplos de aplicación compilan y se han probado con modelos físicos muy sencillos; los nombres de sus sensores y actuadores son de ejemplo.
\item El archivo \texttt{.zal20} de la versión 1.0.0 contiene el error de la cinta y no debe usarse.
\end{itemize}
